% !PRINTED 42
\weberpage{42}

%<<CONT
and we now count again how many distinct elements are contained
in this form. The two values for $z$ must then agree.

If we count, as above in (7), the number of the distinct
elements contained in the form
\weq{12}{B_1^{\beta_1 p_s m} B_2^{\beta_2 p_s m} \ldots B_{s-1}^{\beta_{s-1} p_s m}}
then by the assumption (6) this too yields the value $z$. Hence
in the form
\weq{13}{B_s^{\beta_s p_s m} \ldots B_\mu^{\beta_\mu p_s m}}
only the single element $1$ can be contained, from which it is to be
concluded that $p_s$ must be divisible by $p_s'$. But this is compatible
with (10) only under the supposition that
\weq{14}{p_s = p_s' \text{ holds.}}

With this our Theorem II. is proved.

The prime powers contained in the degree numbers of a basis
of $S$ are thus entirely independent of the particular choice of the
basis, and we therefore call them the \gesperrt{invariants of the
group}. The product of all the invariants is equal to the degree $n$
of the group.

By \S.\,9 we can determine a basis for $S$ in which the
degrees of the elements are all prime powers. By
Theorem II. these degrees are the invariants of the group,
and are thus completely determined by the group.

That the invariants also completely determine the nature
of the group follows from the theorem:

III. \gesperrt{Two groups with the same invariants are
isomorphic, and isomorphic groups have the same
 invariants.}

For if two groups $S$ and $S'$ have basis elements agreeing
in number and in degree, if say
\[ \Theta = A_1^{\alpha_1} A_2^{\alpha_2} \ldots A_\nu^{\alpha_\nu} \]
are the elements of $S$, and
\[ \Theta' = B_1^{\alpha_1} B_2^{\alpha_2} \ldots B_\nu^{\alpha_\nu} \]
the elements of $S'$, where in both cases the $\alpha_1, \alpha_2 \ldots \alpha_\nu$ run through
systems of residues with respect to the moduli $a_1, a_2 \ldots a_\nu$, then we need
only let $\Theta$ and $\Theta'$ correspond to one another when
%CONT>>
